\documentclass[12pt,a4paper]{article}
\usepackage[a4paper,margin=2.2cm]{geometry}
\usepackage{fontspec}
\usepackage{polyglossia}
\setmainlanguage{russian}
\setotherlanguage{english}
\setmainfont{Noto Serif}
\setsansfont{Noto Sans}
\setmonofont{DejaVu Sans Mono}
\usepackage{amsmath,amssymb,mathtools}
\usepackage{array,longtable,booktabs}
\usepackage{enumitem}
\usepackage{hyperref}
\hypersetup{colorlinks=true,linkcolor=black,urlcolor=blue,citecolor=black}
\setlist{nosep,leftmargin=*}
\DeclareMathOperator{\Spin}{Spin}
\DeclareMathOperator{\Tr}{Tr}
\newcommand{\TMI}{\mathcal T_{MI}}
\newcommand{\Adm}{\operatorname{Adm}}
\newcommand{\Struct}{\operatorname{Struct}}
\newcommand{\Dist}{\operatorname{Dist}}
\newcommand{\Cond}{\operatorname{Cond}}
\newcommand{\SM}{\mathrm{SM}}
\newcommand{\GR}{\mathrm{GR}}
\newcommand{\QM}{\mathrm{QM}}
\newcommand{\QG}{\mathrm{QG}}
\newcommand{\BH}{\mathrm{BH}}
\newcommand{\DE}{\mathrm{DE}}
\newcommand{\DM}{\mathrm{DM}}
\newcommand{\dd}{\,d}
\title{TMI-ToE 1.0\\\large Интерфейсно-полевая архитектура объединения фундаментальных взаимодействий}
\author{Теория многообразных интерфейсов}
\date{Версия 1.0}
\begin{document}
\maketitle
\tableofcontents
\newpage

\section*{Аннотация}
\addcontentsline{toc}{section}{Аннотация}

Документ фиксирует модуль \textbf{TMI-ToE 1.0} внутри Теории многообразных интерфейсов. Его задача — не объявить завершённую «теорию всего», а задать минимальную согласованную архитектуру, в которой четыре фундаментальных взаимодействия описываются как проекции единой интерфейсной кривизны.

Короткий статус:
\[
\boxed{\mathrm{TMI\mbox{-}ToE\ 1.0}\neq \text{полная Theory of Everything}}
\]
\begin{center}
\fbox{\parbox{0.88\linewidth}{\centering $\mathrm{TMI\mbox{-}ToE\ 1.0}$ — минимальная интерфейсно-полевая архитектура объединения.}}
\end{center}

Главная идея:
\[
\boxed{
F_a=\pi_a(\mathbb F_I),\qquad a\in\{G,S,W,EM\}
}
\]
где $\mathbb F_I$ — универсальная интерфейсная кривизна, а $F_G,F_S,F_W,F_{EM}$ — гравитационный, сильный, слабый и электромагнитный режимы.

\section{Исходный статус модуля}

После разработки физической ветки ТМИ уже имеются отдельные модули:
\begin{itemize}
  \item \textbf{TMI-QG}: интерфейсная геометрическая декогеренция, величина $\Gamma_I$, проверочная гипотеза и статистическая форма.
  \item \textbf{TMI-Dark}: тёмный сектор как режим интерфейсного поля $\chi_I=\chi_0+\delta\chi$.
  \item \textbf{TMI-BH}: чёрная дыра как односторонний причинно-информационный интерфейс.
  \item \textbf{TMI-ToE}: архитектура объединения этих физических слоёв через универсальную связность, кривизну и интерфейсное поле.
\end{itemize}

TMI-ToE не заменяет проверенные теории. Она требует, чтобы общая теория относительности, Стандартная модель и квантовая механика восстанавливались как предельные случаи.

\section{TMI-ToE-0.9: условия согласованности}

Цель версии 0.9 — закрыть не новые идеи, а технические условия допустимости. Если хотя бы одно из этих условий не выполняется, текущая версия TMI-ToE должна быть пересмотрена.

\subsection{Условие восстановления известных теорий}

\subsubsection*{Предел общей теории относительности}

Если интерфейсный сектор выключен:
\[
\chi_I\to0,\qquad \mathcal L_{mix}\to0,
\]
то должно выполняться:
\[
\boxed{S_{TMI-ToE}\to S_{GR}+S_{SM}}.
\]
Если дополнительно убрать материальные и калибровочные поля:
\[
\mathcal L_{SM}\to0,
\]
то:
\[
\boxed{S_{TMI-ToE}\to S_{GR}}.
\]

\subsubsection*{Предел Стандартной модели}

При фиксированной геометрии:
\[
g_{\mu\nu}=g_{\mu\nu}^{fixed}
\]
и выключенном интерфейсном смешивании:
\[
\chi_I\to0,\qquad \zeta_a,\eta_H,y_I\to0,
\]
должно получаться:
\[
\boxed{S_{TMI-ToE}\to S_{SM}}.
\]
То есть низкоэнергетический калибровочный предел обязан содержать
\[
\boxed{SU(3)\times SU(2)\times U(1)}.
\]

\subsubsection*{Квантовый предел}

При фиксированном фоне и слабом интерфейсном секторе должна восстанавливаться стандартная квантовая динамика:
\[
\boxed{i\hbar\frac{\partial}{\partial t}\Psi=\widehat H\Psi}
\]
или соответствующая квантово-полевая форма на фиксированном фоне.

\subsection{Условие универсальной группы $G_I$}

Минимальное требование:
\[
\boxed{G_I\supset \Spin(1,3)\times SU(3)\times SU(2)\times U(1)}.
\]

Для версии 1.0 безопаснее фиксировать не новую простую группу, а структуру включения:
\[
\boxed{G_I=G_{grav}\ltimes G_{SM}\ltimes G_{\chi}}
\]
где
\[
G_{grav}\supset\Spin(1,3),\qquad G_{SM}=SU(3)\times SU(2)\times U(1).
\]

Смысл: $G_I$ — это группа допустимых интерфейсных преобразований, содержащая гравитационные, калибровочные и интерфейсные симметрии.

\subsection{Условие проекций}

Вводятся проекции алгебр:
\[
\boxed{\pi_a:\mathfrak g_I\to\mathfrak g_a,\qquad a\in\{G,S,W,EM\}},
\]
где $\mathfrak g_I=\mathrm{Lie}(G_I)$.

Тогда физические кривизны задаются как:
\[
\boxed{F_G=\pi_G(\mathbb F_I)},\qquad
\boxed{F_S=\pi_S(\mathbb F_I)},
\]
\[
\boxed{F_W=\pi_W(\mathbb F_I)},\qquad
\boxed{F_{EM}=\pi_{EM}(\mathbb F_I)}.
\]

\subsection{Условие калибровочной инвариантности}

Смешивающий член должен сохранять калибровочную структуру Стандартной модели. Безопасная форма:
\[
\boxed{
\mathcal L_{mix}^{gauge}
=-\frac14\sum_{a=1}^{3}\zeta_a\frac{\chi_I^2}{M_I^2}\Tr(F_{a\mu\nu}F_a^{\mu\nu})
}
\]
допустима, если $\chi_I$ является калибровочным синглетом:
\[
\boxed{\chi_I\sim(1,1,0)}.
\]

\subsection{Условие фермионной согласованности}

Грубый член $\chi_I\bar\Psi\Psi$ до нарушения электрослабой симметрии может быть не калибровочно-инвариантным. Поэтому безопасная форма связи с фермионами:
\[
\boxed{
\mathcal L_{\chi\Psi}
=
\sum_f y_{I,f}\frac{\chi_I}{M_I}\bar\Psi_{L,f}H\Psi_{R,f}+h.c.
}
\]
где $H$ — поле Хиггса, а $f$ — индекс типа фермиона и поколения.

\subsection{Условие хиггсовского сектора}

Совместимый потенциал:
\[
\boxed{
V(H,\chi_I)=
-\mu_H^2H^\dagger H+
\lambda_H(H^\dagger H)^2+
V(\chi_I)+
\eta_H\chi_I^2H^\dagger H
}
\]
где
\[
\boxed{V(\chi_I)=V_0+\frac12m_I^2\chi_I^2+\frac{\beta_I}{4}\chi_I^4}.
\]

Эффективный параметр Хиггса:
\[
\boxed{\mu_{H,eff}^2=\mu_H^2-\eta_H\chi_I^2}.
\]
Следовательно, вакуумное значение может получать интерфейсную поправку:
\[
\boxed{v_H^{eff}=v_H(\chi_I)}.
\]

\subsection{Условие отсутствия аномалий}

Если $\chi_I$ является калибровочным синглетом и не добавляет новых хиральных фермионов, то аномальная структура Стандартной модели не нарушается:
\[
\boxed{\mathcal A_{TMI-ToE}=\mathcal A_{SM}=0}.
\]
Если в будущем добавляются новые фермионы или представления $G_I$, нужно требовать:
\[
\boxed{\mathcal A_{G_I}=0}.
\]

\subsection{EFT-статус}

Текущая версия должна рассматриваться как эффективная интерфейсно-полевая теория:
\[
\boxed{TMI\mbox{-}ToE=\text{effective interface-field theory}}
\]
при энергиях
\[
\boxed{E\ll M_I}.
\]
Операторы вида
\[
\frac{\chi_I^2}{M_I^2}F_{\mu\nu}F^{\mu\nu},\qquad
\frac{\chi_I}{M_I}\bar\Psi_LH\Psi_R
\]
являются эффективными операторами, подавленными масштабом $M_I$.

\subsection{Условие устойчивости}

Потенциал должен быть ограничен снизу. Достаточные условия:
\[
\boxed{\beta_I>0,\qquad \lambda_H>0}.
\]
Для смешанного потенциала нужно исключить нестабильный вакуум. Рабочее условие:
\[
\boxed{\eta_H>-2\sqrt{\lambda_H\beta_I/4}}.
\]

\subsection{Связь с другими физическими модулями ТМИ}

\subsubsection*{Связь с TMI-QG}
\[
\boxed{\chi_I\Rightarrow\Gamma_I},\qquad
\boxed{\Delta\Gamma_I=\kappa_Ix_I}.
\]

\subsubsection*{Связь с TMI-Dark}
\[
\boxed{\chi_I=\chi_0+\delta\chi},\qquad
\boxed{\chi_0\Rightarrow DE,\quad \delta\chi\Rightarrow DM}.
\]

\subsubsection*{Связь с TMI-BH}
На горизонте чёрной дыры должны иметь смысл ограничения:
\[
\boxed{\mathbb A_I|_{\mathcal H_g}},\qquad
\boxed{\mathbb F_I|_{\mathcal H_g}}.
\]

\section{TMI-ToE-1.0: паспорт минимальной завершённой архитектуры}

\subsection{Название}

\[
\boxed{TMI\mbox{-}ToE\ 1.0}
\]

Полное название:
\begin{center}
\fbox{\parbox{0.88\linewidth}{\centering Интерфейсно-полевая архитектура объединения фундаментальных взаимодействий}}
\end{center}

Английский вариант:
\begin{center}
\textit{Interface-Field Architecture for the Unification of Fundamental Interactions.}
\end{center}

\subsection{Статус}

\begin{center}
\fbox{\parbox{0.88\linewidth}{\centering Минимальная завершённая архитектура, но не завершённая физическая Theory of Everything.}}
\end{center}

То есть:
\[
\boxed{TMI\mbox{-}ToE\ 1.0\neq\text{полная теория всего}}.
\]

\subsection{Главный объект}

\[
\boxed{
\mathcal M_I^{ToE}=
(\mathcal M_U,\mathcal E_I,G_I,\mathbb A_I,\mathbb F_I,\Psi,H,\chi_I,\Adm_I,\Struct_I)
}
\]

\begin{longtable}{>{\raggedright\arraybackslash}p{0.23\linewidth}>{\raggedright\arraybackslash}p{0.68\linewidth}}
\toprule
Обозначение & Смысл \\
\midrule
$\mathcal M_U$ & пространство-время \\
$\mathcal E_I$ & универсальное интерфейсное расслоение \\
$G_I$ & универсальная интерфейсная группа \\
$\mathbb A_I$ & универсальная интерфейсная связность \\
$\mathbb F_I$ & универсальная интерфейсная кривизна \\
$\Psi$ & полный квантово-полевой сектор \\
$H$ & поле Хиггса \\
$\chi_I$ & интерфейсное поле \\
$\Adm_I$ & множество допустимых физических переходов \\
$\Struct_I$ & структурированный физический результат \\
\bottomrule
\end{longtable}

\subsection{Универсальное расслоение}

\[
\boxed{\pi_I:\mathcal E_I\to\mathcal M_U}
\]

Слой над точкой:
\[
\boxed{\mathcal F_{I,x}=\pi_I^{-1}(x)}.
\]

Рабочее разложение:
\[
\boxed{
\mathcal F_{I,x}=\mathcal F_G\oplus\mathcal F_{SU(3)}\oplus\mathcal F_{SU(2)}\oplus\mathcal F_{U(1)}\oplus\mathcal F_\Psi\oplus\mathcal F_H\oplus\mathcal F_\chi
}
\]

\subsection{Универсальная связность}

\[
\boxed{\mathbb A_I\in\Omega^1(\mathcal M_U,\mathfrak g_I)}.
\]

Рабочая форма:
\[
\boxed{
\mathbb A_I=
\omega^{ab}J_{ab}+e^aP_a+G^AT_A+W^iT_i+BY+A_\chi\Xi
}
\]
где $A_\chi$ — интерфейсная 1-форма. Если $\chi_I$ берётся как скалярное поле, допустимо:
\[
\boxed{A_\chi=d\chi_I}.
\]

\subsection{Универсальная кривизна}

\[
\boxed{\mathbb F_I=d\mathbb A_I+\mathbb A_I\wedge\mathbb A_I}.
\]

Проекции:
\[
\boxed{F_a=\pi_a(\mathbb F_I),\qquad a\in\{G,S,W,EM\}}.
\]

Главная формула:
\begin{center}
\fbox{\parbox{0.88\linewidth}{\centering Четыре фундаментальных взаимодействия являются проекциями единой интерфейсной кривизны.}}
\end{center}

\subsection{Объединённое действие}

\[
\boxed{S_{TMI-ToE}=S_{GR}+S_{SM}+S_\chi+S_{mix}}.
\]

Развёрнуто:
\[
\boxed{
S_{TMI-ToE}=\int_{\mathcal M_U}\sqrt{-g}\left[
\frac{c^4}{16\pi G}(R-2\Lambda)+\mathcal L_{SM}+\mathcal L_\chi+\mathcal L_{mix}
\right]\dd^4x
}
\]

Интерфейсный лагранжиан:
\[
\boxed{
\mathcal L_\chi=
-\frac12g^{\mu\nu}\nabla_\mu\chi_I\nabla_\nu\chi_I
-V(\chi_I)-\frac12\xi_IR\chi_I^2
}
\]

Смешивающий член:
\[
\boxed{
\mathcal L_{mix}=
-\frac14\sum_{a=1}^{3}\zeta_a\frac{\chi_I^2}{M_I^2}\Tr(F_{a\mu\nu}F_a^{\mu\nu})
+\eta_H\chi_I^2H^\dagger H
+\sum_f y_{I,f}\frac{\chi_I}{M_I}\bar\Psi_{L,f}H\Psi_{R,f}+h.c.
}
\]

\subsection{Главная цепочка}

\[
\boxed{
\mathcal E_I\Rightarrow G_I\Rightarrow\mathbb A_I\Rightarrow\mathbb F_I\Rightarrow\{F_G,F_S,F_W,F_{EM}\}\Rightarrow\Adm_I\Rightarrow\mathcal H_I\Rightarrow\Psi\Rightarrow\Dist_{poss}\Rightarrow\Struct_I
}
\]

Смысл цепочки:
\begin{enumerate}
  \item единая интерфейсная структура задаёт допустимые взаимодействия;
  \item взаимодействия задают физические переходы;
  \item квантовое состояние задаёт распределение возможностей;
  \item интерфейс переводит возможности в структуру.
\end{enumerate}

\section{Проверяемые следствия}

\subsection{Следствие 1: поправки к калибровочным связям}

Из калибровочного смешивания:
\[
-\frac14\zeta_a\frac{\chi_I^2}{M_I^2}\Tr(F_{a\mu\nu}F_a^{\mu\nu})
\]
следует эффективная поправка:
\[
\boxed{
\frac{1}{g_{a,obs}^2(E)}=
\frac{1}{g_{a,SM}^2(E)}+
\zeta_a\frac{\chi_I^2(E)}{M_I^2}
}
\]
и:
\[
\boxed{
\Delta_a(E)=
\frac{1}{g_{a,obs}^2(E)}-
\frac{1}{g_{a,SM}^2(E)}=
\zeta_a\frac{\chi_I^2(E)}{M_I^2}
}
\]

Нулевая гипотеза:
\[
\boxed{H_0:\zeta_a=0}.
\]
Гипотеза TMI-ToE:
\[
\boxed{H_{TMI-ToE}:\exists a,\ \zeta_a\neq0}.
\]

\subsection{Следствие 2: поправки к хиггсовскому сектору}

Через потенциал $V(H,\chi_I)$ получаем:
\[
\boxed{\mu_{H,eff}^2=\mu_H^2-\eta_H\chi_I^2}
\]
и возможную проверочную линию:
\[
\boxed{\Delta v_H=v_H^{obs}-v_H^{SM}=F(\eta_H,\chi_I)}.
\]

\subsection{Следствие 3: связь с интерфейсной декогеренцией}

Если:
\[
\Gamma_I=\lambda_I\|\chi_I\|^2
\]
и:
\[
\Delta_a(E)=\zeta_a\frac{\chi_I^2(E)}{M_I^2},
\]
то возникает мост:
\[
\boxed{\Delta_a(E)=\frac{\zeta_a}{\lambda_I M_I^2}\Gamma_I(E)}.
\]

Это сквозное предсказание: одна и та же интерфейсная интенсивность может проявляться и в декогеренции, и в калибровочных поправках.

\section{Критерии поддержки и ослабления}

\subsection{Критерии поддержки}

Модуль получает поддержку, если:
\begin{enumerate}
  \item корректно восстанавливаются $GR$, $SM$ и $QM$;
  \item смешивающие члены сохраняют калибровочные симметрии;
  \item не появляются неустранимые аномалии;
  \item эффекты подавлены масштабом $M_I$ и не противоречат данным;
  \item наблюдаемые поправки имеют общую зависимость от $\chi_I^2$.
\end{enumerate}

\subsection{Критерии ослабления}

TMI-ToE ослабевает, если:
\begin{enumerate}
  \item не восстанавливается $GR$ или $SM$;
  \item $\mathcal L_{mix}$ нарушает калибровочную инвариантность;
  \item появляются неустранимые аномалии;
  \item эффективные поправки противоречат экспериментальным ограничениям;
  \item нет ни одного отличимого проверяемого следствия.
\end{enumerate}

\section{Финальная фиксация}

\begin{center}
\fbox{\parbox{0.88\linewidth}{\centering $TMI\mbox{-}ToE\ 1.0$ закрыта как минимальная согласованная архитектура.}}
\end{center}

Но:
\begin{center}
\fbox{\parbox{0.88\linewidth}{\centering TMI-ToE 1.0 не является завершённой физической Theory of Everything.}}
\end{center}

Правильная формулировка:
\begin{quote}
TMI-ToE-1.0 — это завершённый минимальный архитектурный модуль объединения, в котором четыре взаимодействия представлены как проекции единой интерфейсной кривизны, а интерфейсное поле $\chi_I$ связывает геометрию, калибровочные поля, квантовое состояние, тёмный сектор, чёрные дыры и процесс структурирования.
\end{quote}

Самая короткая итоговая формула:
\[
\boxed{\mathbb F_I=d\mathbb A_I+\mathbb A_I\wedge\mathbb A_I}
\]
\[
\boxed{F_a=\pi_a(\mathbb F_I),\qquad a\in\{G,S,W,EM\}}
\]
\[
\boxed{S_{TMI-ToE}=S_{GR}+S_{SM}+S_\chi+S_{mix}}
\]
\[
\boxed{\chi_I\Rightarrow\{\Gamma_I,\ DarkSector,\ \Delta_a(E),\ BH\mbox{-}interface\}}
\]

\end{document}
